\SourceLeaf{21}
\par
Facile perspicietur, totam sphaeroidis superficiem sic exhauriri, si $p$ extendatur a $0$ usque ad $180^0$, $q$ vero a $0$ usque ad $360^0$. Porro habebimus
\[
\begin{aligned}
\lambda&=-A\sin p, &\quad\lambda'&=0\\
\mu&=B\cos p.\cos q, &\quad\mu'&=-B\sin p.\sin q\\
\nu&=C\cos p.\sin q, &\quad\nu'&=C\sin p.\cos q
\end{aligned}
\]
\[
\begin{aligned}
\mu\nu'-\nu\mu'&=BC\cos p.\sin p=ABC\sin p.\frac{x}{AA}\\
\nu\lambda'-\lambda\nu'&=AC\sin p^2.\cos q=ABC\sin p.\frac{y}{BB}\\
\lambda\mu'-\mu\lambda'&=AB\sin p^2.\sin q=ABC\sin p.\frac{z}{CC}
\end{aligned}
\]
\par
Hinc quoniam $\sin p$ intra limites, quos hic consideramus, ubique fit quantitas positiva, statuere oportet
\[
\mathrm{d}s=\mathrm{d}p.\mathrm{d}q.ABC.\psi.\sin p
\]
\par
Applicando has formulas ad theorema secundum, fit corporis volumen seu (statuendo densitatem $=1$) massa
\[
=\int\!\int\mathrm{d}p.\mathrm{d}q.ABC.\cos p^2.\sin p
\]
sive integrando primo secundum $q$
\[
=2\pi\int\mathrm{d}p.ABC.\cos p^2.\sin p=\tfrac12\pi ABC\int\mathrm{d}p.(\sin p+\sin3p)
\]
quod integrale a $p=0$ usque ad $p=180^0$ est extendendum. Hinc provenit $\tfrac43\pi ABC$, uti aliunde constat.

\Article{13}
\par
Ad determinandam attractionem, quam sphaerois exercet in punctum quodcunque, si attractio cuiusvis elementi quadrato distantiae a puncto attracto reciproce proportionalis supponitur, habemus $fr=\frac{1}{rr}$, $Fr=-\frac{1}{r}$, $\Phi r=r$. Sit attractio sphaeroidis integri secundum directionem axi coordinatarum $x$ parallelam atque oppositam $=X$, statuaturque $X=ABC\xi$. Erit itaque, per theorema tertium,
\[
X=\int\!\int\mathrm{d}p.\mathrm{d}q\frac{BCx\sin p}{rA}=\int\!\int\mathrm{d}p.\mathrm{d}q\frac{BC\cos p.\sin p}{r}
\]
adeoque
